\documentclass[11pt,a4paper]{article}

\usepackage[T1]{fontenc}
\usepackage[utf8]{inputenc}
\usepackage{lmodern}
\usepackage[margin=2.5cm]{geometry}
\usepackage{microtype}
\usepackage{enumitem}
\usepackage{xcolor}
\usepackage{hyperref}

\setlength{\parindent}{0pt}
\setlength{\parskip}{0.6em}
\setlist[itemize]{topsep=0.2em,itemsep=0.2em}

\definecolor{linkblue}{HTML}{1F4E79}
\hypersetup{
    colorlinks=true,
    linkcolor=linkblue,
    urlcolor=linkblue,
    pdftitle={VIP Weeks 1--2 Report},
    pdfauthor={Yorgo Azzi}
}

\begin{document}

\begin{center}
    {\Large\bfseries VIP Weeks 1--2 Report}\\[0.5em]
    {\large Yorgo Azzi}\\[0.25em]
    September 22, 2026
\end{center}

\section{Overview}

During these two weeks, I focused on adapting my Markov-chain symbolic music
generator for movement-to-sound interaction. The original program generated
music offline and included a command-line realtime mode, but it was not yet
suited to continuous sensor input. The goal was to create a responsive system
in which physical movement could control clear musical properties with minimal
delay.

For the first physical prototype, I used a laptop camera and a colored ribbon
instead of phone sensors. This provides a simple way to test whether the
movement-to-music relationship works before moving to a phone or rope-based
interface.

\section{Work Completed}

\subsection{Realtime Music Engine}

I separated the music-generation logic from the command-line interface and
created a persistent realtime engine. The engine keeps the Markov history,
chord progression, and SATB voice-leading state between musical events instead
of regenerating a complete piece each time a control changes.

I also added direct MIDI playback and a realtime performance runner. Control
values can now change while the music continues playing, without waiting for a
full bar or entering commands between generated sections.

\subsection{Musical Controls and Rhythm}

The original \texttt{intensity} control affected more than one musical
property, so I divided it into two independent controls:

\begin{itemize}
    \item \texttt{volume} controls MIDI loudness;
    \item \texttt{density} controls how fast the rhythm feels.
\end{itemize}

At first, density directly forced fixed note values such as quarter, eighth,
or sixteenth notes. I replaced this with a hybrid approach that keeps the
trained rhythm Markov model active. High density increases the probability of
shorter learned durations, while low density favors longer durations. This
makes the rhythm respond to movement without becoming completely mechanical.

\subsection{Camera Movement Interface}

I added an OpenCV interface that tracks a brightly colored ribbon or band in
the webcam image. The program isolates the selected color, finds the largest
matching blob, and estimates its position and movement speed.

The current mapping is intentionally limited to two controls:

\begin{itemize}
    \item the ribbon's movement speed controls rhythmic density;
    \item the ribbon's vertical position controls volume, with higher positions
    producing louder music.
\end{itemize}

No other movement features currently affect the music. This keeps the first
experiment easy to understand and evaluate.

\section{Results}

The main results were:

\begin{itemize}
    \item The existing music model was fast enough for realtime use. Generating
    a four-beat chorale bar took approximately 40~ms in the initial benchmark.
    \item The new event-based engine generated a musical tick in approximately
    13.8~ms on average, which is comfortably below the available tick time at
    120 beats per minute.
    \item Volume and density can now be changed independently while the music
    is playing.
    \item Faster motion produces shorter notes on average, while the Markov
    model continues to provide rhythmic variation.
    \item The OpenCV tracker successfully detected a synthetic colored blob and
    produced position and motion values.
    \item All 12 automated tests for the realtime engine, rhythm behavior, and
    camera mapping passed.
\end{itemize}

The prototype now connects the complete path from movement tracking to musical
control: camera input, movement analysis, Markov-based generation, and realtime
MIDI output.

\section{Challenges}

Several practical issues still need to be tested:

\begin{itemize}
    \item Camera tracking may change with lighting, background color, and motion
    blur.
    \item End-to-end delay has not yet been measured with a physical webcam and
    audible MIDI output.
    \item The movement-speed sensitivity and smoothing will need adjustment
    through real use.
    \item Realtime performances and movement data are not currently recorded
    for later comparison.
\end{itemize}

\section{Next Steps}

\begin{itemize}
    \item Test the colored ribbon with the laptop webcam under different
    lighting conditions.
    \item Tune the color range, minimum blob size, speed sensitivity, and
    smoothing values.
    \item Measure the delay between a physical gesture and the resulting sound.
    \item Record MIDI output and movement values so different mappings can be
    compared.
    \item Evaluate whether the speed--density and height--volume mappings feel
    natural and produce musically useful results.
    \item Use these results to decide whether to continue with phone sensors or
    begin integrating a rope-based interface.
\end{itemize}

\section{Resources}

\begin{itemize}
    \item Project repository:
    \url{https://github.com/YorgoAzzi/MarkovModel}
    \item OpenCV documentation:
    \url{https://docs.opencv.org/}
\end{itemize}

\end{document}
